% =====================================================================
% PAPER B.  Is the 25 mK anomaly of Ce2Hf2O7 the emergent photon?
% Draft.  Numbers marked \pending{} are NOT yet computed -- see README.md.
% =====================================================================
\documentclass[aps,prb,preprint,amsmath,amssymb,superscriptaddress,
longbibliography,floatfix]{revtex4-2}

\usepackage{bm}
\usepackage{graphicx}
\usepackage{xcolor}
\usepackage[colorlinks=true,citecolor=blue!55!black,
urlcolor=blue!55!black,linkcolor=blue!55!black]{hyperref}

\newcommand{\Jzz}{J_{zz}}
\newcommand{\Jpm}{J_{\pm}}
\newcommand{\Jpmpm}{J_{\pm\pm}}
\newcommand{\CHO}{\mathrm{Ce_2Hf_2O_7}}
\setlength{\emergencystretch}{3em}
\newcommand{\pending}[1]{{\color{red}\textbf{[PENDING: #1]}}}

\begin{document}

\title{Is the 25~mK anomaly of $\mathrm{Ce_2Hf_2O_7}$ the emergent photon?\\
A torus-artifact-free exact-diagonalization test and a parameter-free
entropy budget}

\author{Zhengbang Zhou}
\email{zhengbang.zhou@mail.utoronto.ca}
\affiliation{Department of Physics, University of Toronto, Toronto,
Ontario M5S 1A7, Canada}

\author{Yong Baek Kim}
\email{ybkim@physics.utoronto.ca}
\affiliation{Department of Physics, University of Toronto, Toronto,
Ontario M5S 1A7, Canada}

\date{\today}

\begin{abstract}
The heat capacity of $\CHO$ shows two peaks, at $T_1\simeq64.6$~mK and
$T_2\simeq24.2$~mK, and the lower one has been discussed as the crossover into
the quantum-spin-ice gauge sector.  We test that identification with exact
diagonalization of the nearest-neighbour dipole--octupole XYZ model, after
removing a finite-size effect that has not previously been separated from the
physics.  A periodic cluster is a torus and supports ice-preserving spin
rearrangements that close only through a periodic image; they enter at order
$\Jpm^2/\Jzz$, one order \emph{before} the physical hexagon ring exchange at
$12|\Jpm|^3/\Jzz^2$, and they manufacture the dominant low-temperature scale.
At the published NLC-A couplings this artifact places a spurious heat-capacity
peak at $28.7$~mK --- within $19\%$ of the measured $T_2$, and therefore
capable of confirming the gauge interpretation for the wrong reason.  Deleting
the artifact exactly, by masking transport-changing matrix elements of the
Feshbach map, moves the nearest-neighbour gauge crossover to $6.0$~mK.  Across
$66$ certified points of the fit-constrained coupling region the largest
attainable ratio $T_{\rm gauge}/T_{\rm spinon}$ is $0.198$, against the
measured $0.374$.  We supplement this model-dependent statement with a
model-independent one: any crossover that is the ordering of the ice manifold
must release the Pauling entropy $\tfrac{R}{2}\ln\tfrac32=1.687$~J\,mol$^{-1}$
K$^{-1}$, i.e.\ $29\%$ of $R\ln2$, and any gauge-sector crossover must be
followed below it by a photon $T^3$ tail.  We show what the currently available
digitized data can and cannot decide, and specify the measurement --- heat
capacity below $20$~mK --- that settles it.  Our conclusion is that within the
nearest-neighbour XYZ model the $25$~mK feature is not the gauge crossover; the
honest alternatives are beyond-nearest-neighbour exchange, a larger $|\Jpm|$
than the high-temperature fit constrains, or a non-gauge origin.
\end{abstract}

\maketitle

\section{Introduction}

$\CHO$ is one of a small set of dipole--octupole pyrochlores in which a
quantum-spin-ice (QSI) ground state is plausible on symmetry, chemistry and
spectroscopic grounds~\cite{HuangChenHermele2014,Gao2019,Gaudet2019,%
Bhardwaj2022,RauGingras2019}.  The two-peak heat capacity reported by Smith
\textit{et al.}~\cite{Smith2025} is the sharpest quantitative handle available:
an upper feature near $T_1\simeq65$~mK, naturally read as the spinon (magnetic
monopole) crossover, and a lower feature near $T_2\simeq25$~mK for which the
emergent photon is the natural candidate.  If correct, this is the first
thermodynamic sighting of the QSI gauge sector.

The claim is calibrated by numerics, and the numerics have a problem that is
topological rather than statistical.  Every finite periodic cluster of a
pyrochlore is a three-torus.  On a torus there exist ice-preserving spin
rearrangements whose flipped set closes only after a periodic identification.
Such a process does not exist on the infinite lattice, and on the clusters used
in practice it is \emph{shorter} than the physical process it competes with.
Its amplitude and that of the hexagon ring exchange are
\begin{equation}
 g_4=\frac{4\Jpm^2}{\Jzz},\qquad
 g_6=\frac{12|\Jpm|^3}{\Jzz^2},\qquad
 \frac{g_4}{g_6}=\frac{\Jzz}{3|\Jpm|},
 \label{eq:g4g6}
\end{equation}
so the artifact does not shift the gauge scale, it \emph{replaces} it, with a
different power of the coupling.  At the couplings relevant to $\CHO$,
$|\Jpm|/\Jzz\simeq0.125$, the ratio is $g_4/g_6\simeq2.7$ and the artifact wins.

This paper does three things.  Section~\ref{sec:artifact} shows that at the
published NLC-A couplings the artifact peak lands at $28.7$~mK --- on top of the
measured $24.2$~mK anomaly --- so a small-cluster calculation appears to confirm
the photon interpretation while measuring the box.  Section~\ref{sec:masked}
removes the artifact exactly and reports what the nearest-neighbour model then
predicts.  Section~\ref{sec:entropy} steps outside the model entirely and gives
two parameter-free constraints, an entropy budget and a $T^3$ tail, that the
data must satisfy for \emph{any} gauge interpretation.  Section~\ref{sec:caveat}
states the limitations honestly; they are substantial and they all push in the
same direction.

\section{Model, couplings, and the digitized data}
\label{sec:model}

We use the nearest-neighbour pseudospin Hamiltonian of
Eq.~(\ref{eq:model}) below in the exchange-diagonal local frame,
\begin{equation}
 H=\Jzz\!\!\sum_{\langle ij\rangle}\!S^z_iS^z_j
 -\Jpm\!\!\sum_{\langle ij\rangle}\!\left(S^+_iS^-_j+\mathrm{h.c.}\right)
 +\Jpmpm\!\!\sum_{\langle ij\rangle}\!\left(S^+_iS^+_j+\mathrm{h.c.}\right),
 \label{eq:model}
\end{equation}
related to the published $(J_a,J_b,J_c)$ parametrization by
\begin{equation}
 \Jzz=J_a,\qquad \Jpm=-(J_b+J_c)/4,\qquad \Jpmpm=(J_b-J_c)/4 .
\end{equation}
The three published seeds of Ref.~\cite{Smith2025} are, in meV,
\begin{equation}
 \begin{aligned}
  \text{NLC-A}&=(0.050,\ 0.021,\ 0.004),\\
  \text{NLC-B}&=(0.051,\ 0.008,-0.018),\\
  \text{QMC (ordered)}&=(0.046,-0.003,-0.010).
 \end{aligned}
\end{equation}
NLC-A maps to $(\Jpm,\Jpmpm)/\Jzz=(-0.125,+0.085)$.  Note that $\Jpmpm$ is not
small: omitting it is not a controlled approximation for this material, so all
calculations below use the full XYZ Hamiltonian and its complete $2^{16}$
cubic-16 basis where required.

Measured peak positions, extracted from the published figure, are
$T_2=24.19$~mK and $T_1=64.65$~mK, giving
\begin{equation}
 \mathcal R_{\rm exp}\equiv T_2/T_1=0.374 .
 \label{eq:Rexp}
\end{equation}
The digitized trace begins at $T=20.6$~mK with
$C_{\rm mag}=1.93$~J\,mol$_{\rm Ce}^{-1}$K$^{-1}$, rises to $2.25$ at
$23.7$~mK and falls thereafter.  \textbf{There is no digitized data below the
low-temperature peak}, a fact that limits Sec.~\ref{sec:entropy} and that we do
not paper over.  The extraction carries no experimental uncertainties; it is
used for visual comparison and for the integrals of Sec.~\ref{sec:entropy}, not
for a $\chi^2$.

\section{The artifact lands on the anomaly}
\label{sec:artifact}

Table~\ref{tab:artifact} compares, at fixed couplings and on the same cluster
and the same virtual space, the low-temperature heat-capacity peak of the naive
periodic calculation with the transport-masked one of Ref.~\cite{ZhouKimMask}.
The masking deletes every ice-block matrix element of the exact Feshbach map
$F(z)=PHP+PHQ(z-QHQ)^{-1}QHP$ that connects different transport (polarization)
sectors; a contractible rearrangement preserves the label, a rearrangement
closed through a periodic image changes it by a box vector.  No counterterm and
no fitted coefficient is involved.

\begin{table}[t]
\caption{Low-temperature heat-capacity peak at the $\CHO$-relevant couplings,
in mK, using $\Jzz=J_a$ of the corresponding seed.  ``periodic'' is naive exact
diagonalization of the same cluster and virtual space; ``masked'' has the torus
artifact deleted.  The measured anomaly is at $24.2$~mK.}
\label{tab:artifact}
\begin{ruledtabular}
\begin{tabular}{lccc}
$(\Jpm,\Jpmpm)/\Jzz$ & periodic (mK) & masked (mK) & masked/$g_6$\\ \hline
$(-0.050,\ 0.000)$ & $7.52$  & $0.71$ & $0.821$\\
$(-0.125,\ 0.085)$ (NLC-A) & $17.40$ & $6.01$ & $0.442$\\
$(-0.150,\ 0.000)$ & $17.40$ & $7.72$ & $0.329$\\
$(-0.152,\ 0.085)$ & $17.40$ & $8.11$ & $0.335$\\
\end{tabular}
\end{ruledtabular}
\end{table}

A second, independent extraction on the same cluster with the full two-peak
analysis gives, at NLC-A, a periodic low peak at $28.7$~mK against a masked one
at $6.0$~mK, and at $(-0.152,0.085)$ a periodic $37.5$~mK against a masked
$8.1$~mK.  The two extractions differ in how the peak is located in a
multi-feature curve; the qualitative statement is identical and is the point of
this section:

\begin{quote}
\emph{At the published couplings, naive periodic exact diagonalization predicts
a low-temperature heat-capacity peak within $19\%$ of the measured $24.2$~mK
anomaly.  Essentially all of that peak is the torus artifact.  Removing it moves
the prediction down by a factor of $3$--$5$.}
\end{quote}

The coincidence is not a small perturbation of a correct answer.  Because
$g_4\propto\Jpm^2$ and $g_6\propto|\Jpm|^3$, the artifact carries a different
power of the coupling: a log--log fit of the peak position against $|\Jpm|$
gives $1.80$ for the periodic branch and $2.79$ for the masked branch,
approaching the second and third orders of Eq.~(\ref{eq:g4g6}) with no fitted
subtraction.  A calculation that reproduces $T_2$ through the artifact would
also reproduce the wrong dependence on pressure, chemical substitution, or any
other knob that moves $\Jpm/\Jzz$.

The same artifact has a second, unrelated symptom worth recording for anyone
calibrating on small clusters: it inverts the measured gauge flux.  The winding
amplitude $g_4$ is even in $\Jpm$ while $g_6$ is odd, so for $\Jpm<0$ --- the
sign relevant to every Ce pyrochlore --- artifact and physics oppose and the
artifact wins.  On the $16$-site cluster the raw ground state reports
$\langle O_{\rm hex}\rangle>0$ ($0$-flux) over the entire $(J_x,J_y)$ plane,
whereas the winding-free state reports $\pi$-flux throughout $\Jpm<0$.

\section{What the nearest-neighbour model predicts once the artifact is gone}
\label{sec:masked}

\subsection{The peak-ratio gate}

The absolute position of a finite-cluster gauge peak carries a host-dependent
prefactor: the combinatorial hexagon-flippability coefficient is $1.070$ on the
$16$-site cubic cluster and $0.521$ on the $32$-site FCC cluster, so the
prediction moves down by roughly a factor of two on the larger host and is not
converged even there.  The \emph{ratio} of the two peaks is a better-controlled
observable, because the spinon peak sits at $\simeq0.25\,\Jzz$ and is untouched
by the mask (we verify it moves by less than $1\%$ for $|\Jpm|/\Jzz\le0.09$).

We therefore scan the fit-constrained region and ask whether any certified point
reaches $\mathcal R_{\rm exp}=0.374$.  Of $80$ grid points, $66$ are certified
in the sense that the complete ice-descended band is resolved (all $90$ roots on
cubic-16); the remainder dissolve into the spinon continuum and cannot support
any peak statement at all.  The result is
\begin{equation}
 \max_{\rm certified}\ \mathcal R_{\rm masked}=0.198
 \quad\text{at}\quad (\Jpm,\Jpmpm)/\Jzz=(+0.050,+0.085),
 \label{eq:maxratio}
\end{equation}
with the next two values $0.121$ and $0.111$.  The maximum sits at
\emph{positive} $\Jpm$, i.e.\ in the $0$-flux sector, and is still short of the
measurement by a factor $1.9$.  Restricted to the physically favoured
$\pi$-flux side the shortfall is far larger.

\pending{The repository contains a second, more conservative estimate of the
same quantity --- ``largest reachable ratio $0.050$ at $\Jpm/\Jzz=-0.17$,
shortfall $7.5\times$'' --- computed with a different peak-identification rule.
The two must be reconciled, and one rule adopted, before submission.  Both
support the null; they differ by a factor of four in how badly.}

\subsection{Best available fit}

Ranking the winding-free curves against the digitized $C_{\rm mag}$ over
$0.025$--$2.35$~K by unweighted RMSE in $\log T$ selects
$(\Jpm,\Jpmpm)/\Jzz=(-0.156,+0.115)$ with $\Jzz=0.236$~K $=0.0203$~meV.  That
point puts its two winding-free maxima at $3.53$~mK and $67.2$~mK.  The high
peak is essentially on the measured $T_1$; the low peak is $6.9$ times below
the measured $T_2$.  The nearest-neighbour model can reproduce the spinon scale
and the high-temperature tail while missing the low peak by almost an order of
magnitude.

This is consistent with an independent observation about the fit itself.  At
high temperature the heat capacity depends on the exchange parameters only
through $\sum J^2$ at leading order, so raising $|\Jpm|$ while lowering $J_a$ to
hold $\sum J^2$ fixed is a nearly flat direction of the high-temperature
objective.  An order-$5$ exact numerical linked-cluster calculation reproduces
the published $0.25$--$1.2$~K data with mean residual $2.2\%$ at a point with
$(\Jpm,\Jpmpm)/J_a=(-0.156,0.115)$, better than either published minimum
($3.7\%$ for A, $5.1\%$ for B).  \textbf{The published fit does not constrain
the gauge scale}, because the gauge scale is set by $|\Jpm|^3$ and lives an
order of magnitude below the fitted window.  The material is, on this evidence,
more quantum than the fit says --- but not by enough to close the factor of
seven.

\section{Two constraints that do not depend on the model}
\label{sec:entropy}

Because the model-dependent test above inherits a finite-size prefactor, we add
two constraints that any gauge interpretation must satisfy regardless of
Hamiltonian.

\subsection{The Pauling entropy budget}

If the lower peak is the ordering of the emergent gauge field within the ice
manifold, then integrating the magnetic heat capacity across it must release
exactly the residual ice entropy~\cite{Pauling1935,Ramirez1999},
\begin{equation}
 S_{\rm ice}=\frac{R}{2}\ln\frac32
 =1.687\ \mathrm{J\,mol_{Ce}^{-1}K^{-1}}
 =0.293\,R\ln2 .
 \label{eq:pauling}
\end{equation}
Equivalently, of the total $R\ln2=5.763$~J\,mol$^{-1}$K$^{-1}$ available to an
effective spin-$1/2$ doublet, $71\%$ must be released above the ice plateau (by
the spinon peak and the high-temperature tail) and $29\%$ below it.  This is a
counting statement about the ice manifold, not about any effective gauge
Hamiltonian, and it holds for the $\pi$-flux and $0$-flux sectors alike.

The test is then:
\begin{equation}
 \Delta S_{\rm low}=\int_0^{T_{\rm dip}}\frac{C_{\rm mag}(T)}{T}\,dT
 \ \stackrel{?}{=}\ S_{\rm ice},
 \label{eq:test}
\end{equation}
with $T_{\rm dip}$ the minimum between the two peaks.  A result significantly
below $S_{\rm ice}$ means the low peak is not the ice-manifold crossover
(something else, e.g.\ a fraction of sites or a disorder-broadened subset, is
ordering); a result significantly above it means the peak is not confined to
the ice manifold.

\pending{Eq.~(\ref{eq:test}) has not yet been evaluated.  The obstruction is
real: the digitized trace starts at $20.6$~mK, above the peak's low-temperature
flank, and the unmeasured tail is not negligible --- with
$C_{\rm mag}(20.6~\mathrm{mK})=1.93$, a $T^3$ tail contributes
$\simeq0.64$ and a linear tail $\simeq1.93$~J\,mol$^{-1}$K$^{-1}$, i.e.\
between $38\%$ and $114\%$ of $S_{\rm ice}$ by itself.  The partial integral
from $20.6$~mK to the dip is computable now and is worth reporting with the
tail treated as a stated systematic; a decisive test requires author data below
$20$~mK.  Script: \texttt{analysis/entropy\_budget.py}.}

\subsection{The photon $T^3$ tail}

Below a genuine gauge crossover the specific heat of a linearly dispersing
emergent photon with two polarizations must be
\begin{equation}
 \frac{C}{N}=\frac{2\pi^2}{15}\left(\frac{k_BT}{\hbar c}\right)^3
 \frac{V}{N} ,
 \label{eq:T3}
\end{equation}
with $c$ the emergent light speed~\cite{Benton2012}.  This is the most
distinctive fingerprint the gauge sector has, and it is a two-parameter-free
prediction once $c$ is known: the exponent is $3$ and the coefficient is fixed
by $c$, which is itself computable from the ring-exchange amplitude.  A
measurement of $C_{\rm mag}$ from $5$ to $20$~mK therefore does more than fill
in a tail --- it discriminates the photon ($T^3$) from a Schottky anomaly
(activated), a nuclear contribution ($T^{-2}$), a spin-glass or disorder
contribution (roughly linear), and a conventional magnon ($T^3$ but with a
coefficient set by the ordered moment and a magnetic Bragg peak to accompany
it).  We regard this, not the peak position, as the decisive measurement.

\section{Limitations, and which way they push}
\label{sec:caveat}

\begin{enumerate}
\item \textbf{Cluster size.}  The coupling-plane scan is on the $16$-site
cluster, the only one on which the complete $\Jpmpm$ basis is affordable.  The
host coefficient falls from $1.070$ to $0.521$ between that cluster and the
$32$-site FCC cluster, i.e.\ the larger and more reliable host pushes the gauge
peak \emph{colder}, widening the discrepancy rather than closing it.  This is
the single most important caveat and it does not rescue the gauge reading.
\item \textbf{Virtual-depth truncation.}  On the larger host the virtual space
must be truncated; the measured error of the shallowest truncation is up to
$40\%$ low at strong coupling, again in the direction of a colder peak.
\item \textbf{Band completeness.}  For $\Jpm/\Jzz\lesssim-0.19$ the ice-descended
band is no longer isolated from the spinon continuum on this cluster and no peak
statement is admissible; the minimum ice-manifold overlap falls from $0.89$ at
$\Jpm/\Jzz=-0.05$ to $0.36$ at $-0.18$.  The scan is gated on this, which is why
$14$ of $80$ grid points are excluded.
\item \textbf{Digitized data.}  No experimental uncertainties are available, so
we quote no $\chi^2$ and no confidence interval, and the entropy integral is
incomplete below $20.6$~mK.
\item \textbf{Nearest-neighbour only.}  The model contains no further-neighbour
or dipolar exchange.  The experimental paper itself raises beyond-NN terms as a
possibility, and our result is one of the reasons to take that seriously.
\end{enumerate}

Every numerical caveat pushes the predicted gauge peak to lower temperature.
The null result is therefore robust in the only direction that matters.

\section{Conclusion}

Within the nearest-neighbour dipole--octupole XYZ model, and after the torus
artifact is removed exactly, no certified set of couplings in the
fit-constrained region places a gauge-sector heat-capacity feature at $25$~mK
relative to a $65$~mK spinon feature: the best attainable ratio is $0.198$
against a measured $0.374$, and the best RMSE fit misses the low peak by a
factor of seven.  The coincidence that makes this worth stating carefully is
that the \emph{uncorrected} periodic calculation does place a peak at
$28.7$~mK, so a small-cluster calibration confirms the gauge interpretation by
reproducing its own finite-size artifact.

We do not claim that $\CHO$ is not a quantum spin ice; the spinon scale and the
high-temperature thermodynamics are reproduced well, and the flat direction of
the high-temperature fit means the true $|\Jpm|$ may be appreciably larger than
published.  We claim that the $25$~mK feature is not accounted for as dressed
gauge dynamics of the nearest-neighbour model, and that the question is settled
not by a better fit but by a measurement: the heat capacity between $5$ and
$20$~mK, whose exponent and entropy content decide the matter without reference
to any Hamiltonian.

\begin{acknowledgments}
Heat-capacity data digitized from Ref.~\cite{Smith2025}; the winding-free
construction is that of Ref.~\cite{ZhouKimMask}.
\end{acknowledgments}

\bibliographystyle{apsrev4-2}
\bibliography{refs}

\end{document}
